\documentclass[11pt]{article}

\usepackage[a4paper,margin=1in]{geometry}
\usepackage{amsmath}
\usepackage{amssymb}
\usepackage{hyperref}
\usepackage{graphicx}
\usepackage{listings}
\usepackage{enumitem}

\title{Greedy Graph Coverage for Bus Network Traversal}
\author{Ryan Pek}
\date{\today}

\begin{document}

\maketitle

\section{Problem Definition}

We model the Singapore bus network as a directed graph:

\begin{itemize}
    \item Vertices: bus stops
    \item Directed edges: valid next-stop transitions from bus routes
\end{itemize}

The goal is to construct a traversal that:

\begin{enumerate}
    \item Visits every valid bus stop at least once
    \item Minimizes total travel cost (initially edge count, later time)
\end{enumerate}

This corresponds to a \textbf{graph coverage problem}:

\begin{quote}
Find a walk that visits all vertices at least once, minimizing repetition.
\end{quote}

\section{Baseline Approach}

The simplest approach is greedy:

\begin{itemize}
    \item From current stop, move to the nearest unvisited neighbour
\end{itemize}

This fails due to:

\begin{itemize}
    \item Directed graph constraints
    \item Dead-ends and cycles
    \item Local minima (getting trapped in visited regions)
\end{itemize}

\section{Improved Method: Path-Greedy Traversal}

We define the following algorithm:

\subsection*{Algorithm}

At each step:

\begin{enumerate}
    \item Let $v$ be the current stop
    \item Find the nearest unvisited valid stop $u$ using BFS (or Dijkstra)
    \item Compute shortest path $P(v, u)$
    \item Traverse $P$, allowing revisits
    \item Mark all newly visited stops
    \item Repeat until all valid stops are visited
\end{enumerate}

If no path exists:

\begin{itemize}
    \item Restart at another unvisited stop (new segment)
\end{itemize}

\subsection*{Key Idea}

\begin{quote}
Replace local greedy decisions with global shortest-path jumps.
\end{quote}

\section{Graph Model}

Let:

\begin{itemize}
    \item $G = (V, E)$ be a directed graph
    \item $V$ = set of bus stops
    \item $E$ = directed edges from bus routes
\end{itemize}

Define valid stops:

\[
V_{valid} = \{ v \in V \mid \deg^{+}(v) > 0 \}
\]

\section{Metrics}

We evaluate algorithms using:

\subsection*{Coverage Ratio}

\[
\text{coverage} = \frac{\text{unique visited stops}}{|V_{valid}|}
\]

\subsection*{Efficiency Ratio}

\[
\text{efficiency} = \frac{\text{tour length}}{\text{unique visited stops}}
\]

\subsection*{Revisit Overhead}

\[
\text{overhead} = \frac{\text{tour length} - \text{unique visited stops}}{\text{unique visited stops}}
\]

\section{Experimental Result}

For the Singapore bus graph:

\begin{itemize}
    \item Valid stops: 5017
    \item Tour length: 8334
\end{itemize}

\[
\text{efficiency} = \frac{8334}{5017} \approx 1.66
\]

\[
\text{overhead} \approx 66\%
\]

This indicates:

\begin{itemize}
    \item Full coverage achieved
    \item Moderate repetition due to graph structure constraints
\end{itemize}

\section{Discussion}

The main limitations arise from:

\begin{itemize}
    \item Directed graph structure
    \item Lack of strong connectivity
    \item Local cycles
\end{itemize}

The algorithm avoids local traps by:

\begin{itemize}
    \item Using shortest paths to escape visited regions
    \item Allowing revisits strategically
\end{itemize}

\section{Future Work}

\subsection*{Weighted Graph (Time Optimization)}

Replace BFS with Dijkstra:

\[
\text{edge weight} = \text{travel time}
\]

Possible components:

\begin{itemize}
    \item Bus travel duration
    \item Waiting time
    \item Transfer penalties
\end{itemize}

\subsection*{Graph Decomposition}

Partition graph into clusters:

\begin{itemize}
    \item Strongly Connected Components (SCCs)
    \item Weakly Connected Components
    \item Community clusters
\end{itemize}

Strategy:

\begin{enumerate}
    \item Fully cover local cluster
    \item Move between clusters via high-connectivity nodes
\end{enumerate}

\subsection*{Centrality-Based Optimization}

Use:

\begin{itemize}
    \item Betweenness centrality to identify hubs
    \item Degree heuristics to guide traversal
\end{itemize}

\subsection*{Closed Tour Extension}

To return to start:

\begin{itemize}
    \item Append shortest path from final node back to start
\end{itemize}

This becomes:

\begin{quote}
Directed graph traveling salesman variant (asymmetric TSP)
\end{quote}

\section{Conclusion}

We developed a practical graph traversal method:

\begin{itemize}
    \item Greedy local strategies are insufficient
    \item Global shortest-path guidance is necessary
    \item Directed constraints introduce unavoidable overhead
\end{itemize}

The resulting algorithm:

\begin{quote}
Achieves full coverage with manageable repetition, and serves as a baseline for further optimization toward real-world travel efficiency.
\end{quote}

\end{document}
